% Principal: GDZ PPN236018647, 1917, PDF native 12, printed 5.
\SourceLeaf{12}\NoSourceHead
\pdfbookmark[1]{Vorrede zu Piazzis Lehrbuch der Astronomie}{piazzi}
\HistoricTitle{VORREDE.}
\HistoricCitation{Lehrbuch der Astronomie von \textsc{Joseph Piazzi}.\\
Aus dem Italienischen übersetzt von \textsc{Johann Heinrich Westphal}.\\
Berlin bei G. Reimer 1822.}

Das Original dieses Werks ist im Jahr 1817 unter dem Titel \textls[100]{Lezioni elementari di astronomia ad usu del real osservatorio di Palermo} in zwei Bänden zu Palermo erschienen, und war zunächst für die astronomischen Vorlesungen bestimmt, welche der Verf. zu halten hatte. Obgleich wir an elementarischen Schriften über Astronomie keinen Mangel haben, und einige darunter in ihrer Art vortrefflich sind, so wird man doch, bei der Verschiedenheit der Vorkenntnisse und Absichten der Leser, das Hinzukommen einer neuen nicht für überflüssig halten, zumal von einem Verfasser, der sich um mehrere Theile der Wissenschaft so hohe Verdienste erworben hat. Ich zweifle daher nicht, dass die Übertragung dieses Werks in unsere Sprache, durch einen Gelehrten, der seine gründlichen Einsichten bereits durch eigne Arbeiten erprobt hat, allen Freunden der Astronomie willkommen sein werde, die entweder durch die Sprache, oder durch die Schwierigkeit, sich italienische Werke zu verschaffen, abgehalten werden das Original zu lesen. Hin und wieder hat der Herr Übersetzer in kleinen Einschaltungen einiges hinzugesetzt, was dem Verfasser bei der Herausgabe des Originals noch nicht bekannt sein konnte, und der letzte Abschnitt, die Berechnung der Cometenbahnen betreffend, rührt von jenem allein her, welcher den Lesern einen Dienst zu erzeigen glaubte, wenn er die \textsc{Olbers}sche Methode an die Stelle der von dem Verfasser gewählten unvollkommenen indirecten setzte.

Möge diese Arbeit dazu beitragen, die mehr als oberflächliche Befreundung mit einer Wissenschaft zu befördern, die so vielfachen Stoff zu einer edeln und kräftigen Geistesnahrung darbietet.

\bigskip\noindent Göttingen, den 20. October 1821.\hfill C. F. \textsc{Gauss}.
\clearpage
